% !TEX root = ../Projektdokumentation.tex
\section{Abnahmephase} 
\label{sec:Abnahmephase}
Um zu versichern, dass Logik und Systemkomponenten wie gefordert funktionieren wurden der Applikation
\textit{Nightcrawler} Unit- sowie Integrationstests hinzugefügt. Welche ebenfalls bei jedem push in einen
pull-request durch pipelines ausgeführt werden, um nach jeder Änderung zu gewährleisten, dass die Logik sich nicht
geändert hat. Beim Framework zum Testen handelt es sich um \textit{NSubstitute} mit \textit{xUnit} und der Erweiterung \textit{FluentAssertions},
was ermöglicht Tests in lesbarer Form im Arrange, Act und Assert Schema zu schreiben. Außerdem wurden die Tests immer nach dem Muster 
'zu testende public Methode\_Zustand\_Resultat' benannt, um damit einheitliche und verständliche Tests zu gewährleisten.


\subsection{Unittests}
\label{sec:Unittests}
Pro .NET Projekt wurde ein neues Testprojekt für Unittests erstellt.
Diese Art von Tests beschränken sich auf das Überprüfen von einzelnen Codeeinheiten. In diesem Fall wurde eine Unittestklasse
pro zu testender Klasse erstellt. So erhält jede Testklasse eine echte Instanz der zu testenden Klasse, die aber mit 
substituierten/gemockten Abhängigkeiten instanziiert werden, sodass sich beim Testen rein auf die zu testende Logik beschränkt werden konnte. 


\subsection{Integrationstests}
\label{sec:Integrationstests}
Integrationstest bilden den realen Ablauf von Programmen ab. Um echte Kommunikationswege testen zu können wurden die Pakete \textit{Testcontainers.PostgreSql} und 
\textit{Testcontainers.RabbitMQ} verwendet, die es ermöglichen Instanzen der jeweiligen Technologie mit Docker zu starten, welche in den Integrationstest, wie eine 
echte Instanz angesprochen werden können.  
\par\smallskip
Im Anhang finden sich die Projektstruktur und Testbeispiele \Anhang{app:Test}.


\subsection{Test Ergebnisse}
\label{sec:Ergebnisse}
Das Ergebnis dieser Tests ist das Erzielen einer Testabdeckung von 95 Prozent und einer fehlerfreien Anwendung.
\begin{figure}[htb]
\centering
\includegraphics[scale=0.56]{coverage.png}\\
\Anhang{app:Test}
\end{figure} 
